\subsection{Definition of rings of fractions}

DEFINITION 1. Let A be a ring. A subset S of A is called multiplicative if every finite product of elements of S belongs to S.

This is the same as saying that \(1 \in S\) and that the product of two elements of \(S\) belong to \(S\) .
% semantic-fragments: legacy-input-seam
\relax{}
\textbf{Examples}\par

\begin{enumerate}
\def\labelenumi{(\arabic{enumi})}
\item
  For every \(a \in \mathbf{A}\) , the set of \(a^n\) , where \(n \in \mathbf{N}\) , is a multiplicative subset of \(\mathbf{A}\) .
\item
  Let \(\mathfrak{p}\) be an ideal of \(\mathbf{A}\) . For \(\mathbf{A} - \mathfrak{p}\) to be a multiplicative subset of \(\mathbf{A}\) , it is necessary and sufficient that \(\mathfrak{p}\) be prime.
\item
  The set of elements of A which are not divisors of zero is a multiplicative subset of A.
\item
  If S and T are multiplicative subsets of A, the set ST of products st, where \(s \in S\) and \(t \in T\) , is a multiplicative subset.
\item
  Let \(\mathfrak{S}\) be a directed set (with respect to the relation \(\subset\) ) of multiplicative subsets of A. Then \(T = \bigcup_{S \in \mathfrak{S}} S\) is a multiplicative subset of A, as any two elements of T belong to some subset \(S \in \mathfrak{S}\) , hence their product belongs to T.
\item
  Every intersection of multiplicative subsets of A is a multiplicative subset.
\end{enumerate}

For every subset S of a ring A, there exist multiplicative subsets of A containing S, for example A itself. The intersection of all these subsets is the smallest multiplicative subset of A containing S; it is said to be generated by S. It follows immediately that it is the set consisting of all the finite products of elements of S.

PROPOSITION 1. Let A be a ring and S a subset of A. There exists a ring A' and a homomorphism h of A to A' with the following properties:

\begin{enumerate}
\def\labelenumi{(\arabic{enumi})}
\item
  the elements of \(h(\mathbf{S})\) are invertible in \(\mathbf{A}'\) ;
\item
  for every homomorphism \(u\) of \(\mathbf{A}\) to a ring \(\mathbf{B}\) such that the elements of \(u(\mathbf{S})\) are invertible in \(\mathbf{B}\) , there exists a unique homomorphism \(u'\) of \(\mathbf{A}'\) to \(\mathbf{B}\) such that \(u = u' \circ h\) .
\end{enumerate}

In other words, \((A', h)\) is a solution of the universal mapping problem (Set Theory, Chapter IV, §3, no. 1) with the following conditions: the species of structure \(\Sigma\) considered is that of a ring, the morphisms are ring homomorphisms and the \(\alpha\) -mappings are homomorphisms of A to a ring such that the image of S under such a homomorphism consists of invertible elements. Recall (loc. cit.) that, if \((A', h)\) and \((A_{1}', h_{1})\) are both solutions of this problem, there exists a unique isomorphism \(j: A' \to A_{1}'\) such that \(h_{1} = j \circ h\) .

Let \(\bar{S}\) be the multiplicative subset of A generated by S. Clearly every solution of the above universal mapping problem is also a solution of the universal mapping problem obtained by replacing S by \(\bar{S}\) and conversely.

Consider, in the set \(\mathbf{A} \times \overline{\mathbf{S}}\) , the following relation between elements \((a, s)\) , \((a', s')\) :

\begin{enumerate}
\def\labelenumi{(\arabic{enumi})}
\tightlist
\item
 ``There exists \(t \in \overline{\mathbf{S}}\) such that \(t(sa' - s'a) = 0\)''.
\end{enumerate}

This relation is an equivalence relation: it is clearly reflexive and symmetric; it is transitive, for if \(t(sa' - s'a) = 0\) and \(t'(s'a'' - s''a') = 0\) , then \(tt's'(sa' - s'a) = 0\) and \(tt's' \in \overline{\mathbb{S}}\) . Let \(A'\) be the quotient set of \(A \times \overline{S}\) under this equivalence relation; for every ordered pair \((a, s) \in A \times \overline{S}\) we denote by \(a/s\) the canonical image of \((a, s)\) in \(A'\) and we set \(h(a) = a/1\) for all \(a \in A\) . We shall see that \(A'\) can be given a ring structure such that the ordered pair \((A', h)\) solves the problem.

Let \(x = a / s\) and \(y = b / t\) be two elements of \(A'\) . The elements \((ta + sb) / st\) and \(ab / st\) depend only on \(x\) and \(y\) ; for if \(x = a' / s'\) , there exists by hypothesis \(r \in \overline{\mathbb{S}}\) such that \(r(s'a - sa') = 0\) , whence

\[
r \left(s ^ {\prime} t (t a + s b) - s t \left(t a ^ {\prime} + s ^ {\prime} b\right)\right) = 0
\]

and \(r(s'tab - sta'b) = 0\) . It is easily verified that the laws of composition \((x,y) \mapsto x + y = (ta + sb)/st\) and \((x,y) \mapsto xy = ab/st\) define a commutative ring structure on \(A'\) , under which 0/1 is the identity element under addition and 1/1 is the unit element. Moreover, it is immediate that \(h\) is a ring homomorphism and that, for all \(s \in S\) , \(s/1\) is invertible in \(A'\) , its inverse being 1/s. Finally let B be a ring and \(u: A \to B\) a homomorphism such that the elements \(u(S)\) are invertible in B; there exists a unique mapping \(u': A' \to B\) such that

\[
u ^ {\prime} (a / s) = u (a) (u (s)) ^ {- 1} \quad (a \in A, s \in \bar {S}).\tag{2}
\]

If \(a / s = a' / s'\) , there exists \(t \in \overline{S}\) such that \(t(sa' - s'a) = 0\) , whence \(u(t)(u(s)u(a') - u(s'))u(a)) = 0\) and, as \(u(t), u(s)\) and \(u(s')\) are invertible, \(u(a)(u(s))^{-1} = u(a')(u(s'))^{-1}\) . It is easily verified that \(u'\) is a homomorphism with respect to addition and multiplication; finally, clearly \(u' \circ h = u\) and \(u'\) is the only homomorphism satisfying this relation, as it implies

\[
u ^ {\prime} (a / s) = u ^ {\prime} ((a / 1) (1 / s)) = u ^ {\prime} (1 / s) u ^ {\prime} (a / 1) = u ^ {\prime} (1 / s) u (a)
\]

and \(1 = u'(1/1) = u'(s/1)u'(1/s) = u(s)u'(1/s)\) , whence formula (2).

DEFINITION 2. Let A be a ring, S a subset of A and \(\overline{\mathbf{S}}\) the multiplicative subset generated by S. The ring of fractions of A defined by S and denoted by \(\mathrm{A}[\mathrm{S}^{-1}]\) is the quotient set of \(\mathrm{A} \times \overline{\mathrm{S}}\) under the equivalence relation (1) with the ring structure defined by

\[
(a / s) + (b / t) = (t a + s b) / s t, \quad (a / s) (b / t) = (a b) / (s t)
\]

for \(a, b\) in A, \(s, t\) in \(\bar{S}\) . The canonical mapping of A to A{[}S \(^{-1}\) {]} is the homomorphism \(a \mapsto a/1\) , which makes A{[}S \(^{-1}\) {]} into an A-algebra.

In this chapter we usually denote this canonical mapping by \(i_{A}^{S}\) ; the proof of Proposition 1 shows that the ordered pair \((\mathrm{A}[\mathrm{S}^{-1}], i_{\mathrm{A}}^{\mathrm{S}})\) satisfies the conditions of the statement of this proposition.

Remarks

\begin{enumerate}
\def\labelenumi{(\arabic{enumi})}
\item
  Clearly \(\mathbf{A}[\overline{\mathbf{S}}^{-1}] = \mathbf{A}[\mathbf{S}^{-1}]\) .
\item
  Two elements of A{[}S \(^{-1}\) {]} can always be written in the form a/s and a'/s (a, a' in A, s ∈ S) with the same ``denominator'' s, for if b/t and b'/t' are two elements of A{[}S \(^{-1}\) {]}, then b/t = bt'/tt' and b'/t' = b't/tt'.
\item
  The kernel of \(i_{\mathbf{A}}^{\mathbf{s}}\) is the set of \(a \in \mathbf{A}\) such that there exists \(s \in \overline{\mathbf{S}}\) satisfying \(sa = 0\) ; for \(i_{\mathbf{A}}^{\mathbf{s}}\) to be injective, it is necessary and sufficient that \(\mathbf{S}\) contain no divisor of zero in \(\mathbf{A}\) .
\item
  If \(S\) contains a nilpotent element, then \(0 \in \overline{S}\) and the ring \(A[S^{-1}]\) is reduced to 0; this follows easily from Definition 2.
\item
  For \(i_{\mathbf{A}}^{\mathbf{S}}\) to be a bijection, it is necessary and sufficient that every element \(s \in \mathbf{S}\) be invertible in A: the condition is obviously necessary, since \(s/1\) is invertible in \(\mathrm{A}[\mathrm{S}^{-1}]\) ; it is sufficient, since for all \(t \in \overline{\mathrm{S}}\) , \(t\) is therefore invertible in A and \(a/t = at^{-1}/1\) in \(\mathrm{A}[\mathrm{S}^{-1}]\) ; hence \(i_{\mathbf{A}}^{\mathbf{S}}\) is surjective and it has already been seen in Remark 3 that it is injective. Then A and \(\mathrm{A}[\mathrm{S}^{-1}]\) are identified by means of \(i_{\mathbf{A}}^{\mathbf{S}}\) .
\end{enumerate}

Example (7). If R is the set of elements in A which are not divisors of 0, the ring \(\mathrm{A}[\mathrm{R}^{-1}]\) is precisely what we have called the ring of fractions of A (Algebra, Chapter I, §9, no. 4); to avoid any confusion we shall often call it the total ring of fractions of A. In particular, if A is an integral domain, \(\mathrm{A}[\mathrm{R}^{-1}]\) is the field of fractions of A (loc. cit.).

PROPOSITION 2. Let A, B be two rings, S a subset of A, T a subset of B and f a homomorphism from A to B such that \(f(S) \subset T\) . There exists a unique homomorphism \(f'\) from A{[}S \(^{-1}\) {]} to B{[}T \(^{-1}\) {]} such that \(f'(a/1) = f(a)/1\) for all \(a \in A\) .

Suppose further that \(\mathbf{T}\) is contained in the multiplicative subset of \(\mathbf{B}\) generated by \(f(\mathbf{S})\) . Then, if \(f\) is surjective (resp. injective) so is \(f'\) .

The first assertion amounts to saying that there exists a unique homomorphism \(f'\) : \(A[S^{-1}] \to B[T^{-1}]\) giving a commutative diagram:

\[
\begin{array}{c} \mathrm{A} \xrightarrow {f} \mathrm{B} \\ i _ {\mathbf {A}} ^ {\mathbf {S}} \Biggl \downarrow \qquad \qquad \qquad \qquad \Biggl \downarrow i _ {\mathbf {B}} ^ {\mathbf {T}} \\ \mathrm{A} [ \mathrm{S} ^ {- 1} ] \xrightarrow [ f ^ {\prime} ]{} \mathrm{B} [ \mathrm{T} ^ {- 1} ] \end{array}
\]

Now the relation \(f(\mathbf{S}) \subset \mathbf{T}\) implies that \(i_{\mathbf{B}}^{\mathbf{T}}(f(s))\) is invertible in \(\mathbf{B}[\mathbf{T}^{-1}]\) for all \(s \in \mathbf{S}\) and it is sufficient to apply Proposition 1 to \(i_{\mathbf{B}}^{\mathbf{T}} \circ f\) . It follows easily from (2) that, for all \(a \in \mathbf{A}\) and \(s \in \overline{\mathbf{S}}\) (multiplicative subset of \(\mathbf{A}\) generated by \(\mathbf{S}\) ),

\[
f ^ {\prime} (a / s) = f (a) / f (s).\tag{3}
\]

Suppose that \(\mathbf{T}\) is contained in the multiplicative subset generated by \(f(\mathbb{S})\) , which is precisely \(f(\overline{\mathbb{S}})\) . Then it follows from (3) that, if \(f\) is surjective, so is \(f'\) .

Suppose now that \(f\) is injective. Let \(a / s\) be an element of the kernel of \(f'\) . As the multiplicative subset generated by \(T\) is \(f(\overline{S})\) , there is an element \(s_1 \in \overline{S}\) such that \(f(s_1)f(a) = 0\) , whence \(f(s_1a) = 0\) and consequently \(s_1a = 0\) since \(f\) is injective; then \(a / s = 0\) , which proves that \(f'\) is injective.

Remark (6). If the elements of T are invertible in B, \(B[T^{-1}]\) is identified with B by means of the isomorphism \(i_{B}^{T}\) and \(f'\) then becomes identical with the unique homomorphism \(u'\) of \(A[S^{-1}]\) to B such that \(u' \circ i_{A}^{S} = f\) .

COROLLARY 1. Let A be a ring, S a subset of A and u an injective homomorphism of A to a ring B such that the elements of \(u(S)\) are invertible in B. The unique homomorphism \(u'\) of \(\mathbf{A}[\mathbf{S}^{-1}]\) to B such that \(u' \circ i_{\mathbf{A}}^{\mathbf{S}} = u\) is then injective.

This is an immediate consequence of Proposition 2 and Remark 6.

COROLLARY 2. Let A be a ring and S and T two subsets of A such that \(\mathbf{S} \subset \mathbf{T}\) . There exists a unique homomorphism \(i_{\mathbf{A}}^{\mathrm{T},\mathrm{S}}\) from \(\mathrm{A}[\mathrm{S}^{-1}]\) to \(\mathrm{A}[\mathrm{T}^{-1}]\) such that \(i_{\mathbf{A}}^{\mathrm{T}} = i_{\mathbf{A}}^{\mathrm{T},\mathrm{S}} \circ i_{\mathbf{A}}^{\mathrm{S}}\) .

For all \(a \in \mathbf{A}\) , \(i_{\mathbf{A}}^{\mathbf{T},\mathbf{S}}\) then maps the element \(a / s\) in \(\mathbf{A}[\mathbf{S}^{-1}]\) to the element \(a / s\) in \(\mathbf{A}[\mathbf{T}^{-1}]\) .

Remark (7). Note that, if \(i_{A}^{T}\) is injective, so is \(i_{A}^{T,S}\) (Corollary 1). This is what happens if T is the set R of elements of A which are not divisors of 0; then it is possible to identify \(A[S^{-1}]\) with the subring of the total ring of fractions \(A[R^{-1}]\) generated by A and the inverses in \(A[R^{-1}]\) of the elements of S.

COROLLARY 3. Let A, B, C be three rings, S (resp. T, U) a multiplicative subset of A (resp. B, C), \(f\colon A \to B\) , \(g\colon B \to C\) two homomorphisms and \(h\colon A \to C\) the composite homomorphism \(g \circ f\) ; suppose that \(f(S) \subset T\) , \(g(T) \subset U\) . Let \(f'\colon A[S^{-1}] \to B[T^{-1}]\) , \(g'\colon B[T^{-1}] \to C[U^{-1}]\) , \(h'\colon A[S^{-1}] \to C[U^{-1}]\) the homomorphisms corresponding to \(f\) , \(g\) , \(h\) ; then \(h' = g' \circ f'\) .

This follows easily from the definitions.

In particular, if S, T, U are three multiplicative subsets of A such that \(S \subset T \subset U\) , then \(i_{A}^{U,S} = i_{A}^{U,T} \circ i_{A}^{T,S}\) .

COROLLARY 4. Let S be a subset of a ring A, B a subring of A{[}S \(^{-1}\) {]} containing \(i_{\text{A}}^{\text{S}}(\text{A})\) and S' the set \(i_{\text{A}}^{\text{S}}(\text{A})\) . Let j be the canonical injection of B into A{[}S \(^{-1}\) {]}; the unique homomorphism g from B{[}S \(^{-1}\) {]} to A{[}S \(^{-1}\) {]} such that \(g \circ i_{\text{A}}^{\text{S}} = j\) is an isomorphism.

The mapping g is injective by Corollary 1; the ring \(g(\mathbf{B}[\mathbf{S}^{\prime-1}])\) contains \(i_{\mathbf{A}}^{\mathbf{s}}(\mathbf{A})\) and the inverse of the elements of \(S'\) ; hence it is equal to \(A[S^{-1}]\) .

If A is an integral domain and \(0 \notin S\) , the notation \(A[S^{-1}]\) agrees with that of Algebra, Chapter IV, § 2, no. 1; also, if S is multiplicative, \(A[S^{-1}]\) coincides in this case with the set denoted by \(S^{-1}A\) in Algebra, Chapter I, § 1, no. 1.

As an extension of notation, for any multiplicative subset S of a ring A, we henceforth denote by \(S^{-1}A\) the ring of fractions \(A[S^{-1}]\) . If S is the complement of a prime ideal p of A, we write \(A_{p}\) instead of \(S^{-1}A\) .

If A is an integral domain and \(0 \notin S\) , \(S^{-1}A\) is always identified with a subring of the field of fractions of A, containing A (Remark 7).

